% --- CORRECCIÓN IMPORTANTE EN LA PRIMERA LÍNEA ---
% Agregamos "xcolor=table" para que funcione \rowcolor sin errores
\documentclass[aspectratio=169, 10pt, xcolor=table]{beamer}

% --- Configuración del Tema (estilo profesional como el ejemplo) ---
\usetheme{Madrid}
\usecolortheme{dolphin}

% --- Paquetes necesarios ---
\usepackage[utf8]{inputenc}
\usepackage[spanish]{babel}
\usepackage{amsmath, amssymb}
\usepackage{booktabs}
\usepackage{graphicx}
\usepackage{listings}
\usepackage{tikz}
\usetikzlibrary{arrows.meta, positioning, shapes.geometric}

% --- Ruta de Imágenes (crea la carpeta "imagenes" y coloca ahí tus archivos) ---
\graphicspath{{./imagenes/}}

% --- Estilo para código Python ---
\definecolor{codegreen}{rgb}{0,0.6,0}
\definecolor{codegray}{rgb}{0.5,0.5,0.5}
\definecolor{codepurple}{rgb}{0.58,0,0.82}
\definecolor{backcolour}{rgb}{0.95,0.95,0.92}
\lstset{
    language=Python,
    backgroundcolor=\color{backcolour},
    commentstyle=\color{codegreen},
    keywordstyle=\color{blue},
    numberstyle=\tiny\color{codegray},
    stringstyle=\color{codepurple},
    basicstyle=\ttfamily\scriptsize,
    breaklines=true,
    frame=single,
    showstringspaces=false
}

% --- Pie de página personalizado ---
\makeatletter

\setbeamertemplate{footline}{
  \leavevmode%
  \hbox{%
  \begin{beamercolorbox}[wd=.333333\paperwidth,ht=2.25ex,dp=1ex,center]{author in head/foot}%
    \usebeamerfont{author in head/foot}\insertshortauthor
  \end{beamercolorbox}%
  \begin{beamercolorbox}[wd=.333333\paperwidth,ht=2.25ex,dp=1ex,center]{title in head/foot}%
    \usebeamerfont{title in head/foot}Sesión 10
  \end{beamercolorbox}%
  \begin{beamercolorbox}[wd=.333333\paperwidth,ht=2.25ex,dp=1ex,right]{date in head/foot}%
    \usebeamerfont{date in head/foot}NLP \hfill \insertframenumber/\inserttotalframenumber \hspace*{0.3cm}
  \end{beamercolorbox}}%
  \vskip0pt%
}

\makeatother

% --- Datos de la portada ---
\title[SESIÓN 10 NLP]{\textbf{Sesión 10}}
\subtitle{Modelos Pre-entrenados (BERT, GPT) y Fine-tuning}
\author[Prof. C. Sánchez]{Carlos César Sánchez Coronel}
\institute{\textbf{NLP}}
\date{}

% Logo opcional (descomentar si tienes la imagen)
% \titlegraphic{\includegraphics[height=1.5cm]{logo.png}}

% --- Eliminar la barra de navegación (opcional) ---
\setbeamertemplate{navigation symbols}{}

% --- Índice al inicio de cada sección ---
\AtBeginSection[]{
  \begin{frame}
    \frametitle{Contenido}
    \tableofcontents[currentsection]
  \end{frame}
}

% ============================================================================
% INICIO DEL DOCUMENTO
% ============================================================================
\begin{document}

% --- Portada ---
\begin{frame}
    \titlepage
\end{frame}

% --- Índice general ---
\begin{frame}{Contenido}
    \tableofcontents
\end{frame}

% ============================================================================
% SECCIÓN 1: INTRODUCCIÓN AL TRANSFER LEARNING
% ============================================================================
\section{Introducción al Transfer Learning}

\begin{frame}{El poder del pre-entrenamiento}
    \begin{itemize}
        \item Antes de 2018: modelos entrenados desde cero para cada tarea $\rightarrow$ necesitaban muchos datos etiquetados.
        \item Modelos pre-entrenados (BERT, GPT) aprenden representaciones del lenguaje a partir de grandes corpus no etiquetados.
        \item \textbf{Transfer learning}: adaptar (fine-tune) estos modelos a tareas específicas con pocos datos etiquetados.
        \item Es el estándar actual en NLP.
    \end{itemize}
    % Basado en introducción del PDF
\end{frame}

% ============================================================================
% SECCIÓN 2: BERT
% ============================================================================
\section{BERT (Bidirectional Encoder Representations from Transformers)}

\begin{frame}{Arquitectura de BERT}
    \begin{itemize}
        \item Basado en el \textbf{encoder} del Transformer.
        \item \textbf{Bidireccional}: cada token puede atender a tokens a izquierda y derecha (sin máscara causal).
        \item Profundamente bidireccional, a diferencia de modelos anteriores.
    \end{itemize}
    % Basado en PDF y Pregunta 1
    % IMAGEN: Diagrama de BERT
\end{frame}

\begin{frame}{Tokenización: WordPiece}
    \begin{itemize}
        \item Vocabulario de subpalabras (aprox. 30,000 tokens).
        \item Algoritmo: fusiona pares que maximizan la probabilidad del lenguaje:
        \[
        \text{score}(a,b) = \frac{\text{freq}(ab)}{\text{freq}(a) \cdot \text{freq}(b)}
        \]
        \item Maneja palabras desconocidas descomponiéndolas en subpalabras conocidas.
    \end{itemize}
    % Basado en PDF y Pregunta 5
\end{frame}

\begin{frame}{Objetivos de pre-entrenamiento de BERT}
    \begin{columns}[t]
        \column{0.5\textwidth}
        \textbf{Masked Language Model (MLM)}
        \begin{itemize}
            \item Se enmascara el 15\% de los tokens.
            \item El modelo debe predecir los tokens enmascarados.
            \item Fuerza a aprender representaciones contextuales.
        \end{itemize}
        \column{0.5\textwidth}
        \textbf{Next Sentence Prediction (NSP)}
        \begin{itemize}
            \item Predice si dos oraciones son consecutivas.
            \item Útil para tareas que requieren relaciones entre oraciones (QA, inferencia).
        \end{itemize}
    \end{columns}
    \vspace{0.3cm}
    \begin{exampleblock}{Detalles del MLM}
        De los tokens enmascarados: 80\% [MASK], 10\% palabra aleatoria, 10\% se deja igual.
    \end{exampleblock}
    % Basado en PDF y Preguntas 2, 7
\end{frame}

\begin{frame}{Representaciones de salida}
    \begin{itemize}
        \item Para cada token, BERT produce un vector contextualizado.
        \item El token \texttt{[CLS]} (primera posición) se usa como representación de toda la secuencia para tareas de clasificación.
        \item Para tareas de etiquetado (NER, POS), se usan las representaciones de cada token.
    \end{itemize}
    % Basado en PDF y Pregunta 6
\end{frame}

% ============================================================================
% SECCIÓN 3: GPT
% ============================================================================
\section{GPT (Generative Pre-trained Transformer)}

\begin{frame}{Arquitectura de GPT}
    \begin{itemize}
        \item Basado en el \textbf{decoder} del Transformer.
        \item \textbf{Unidireccional (causal)}: cada token solo puede atender a tokens anteriores (máscara causal).
        \item Diseñado para generación de texto.
    \end{itemize}
    % Basado en PDF y Pregunta 1
\end{frame}

\begin{frame}{Tokenización: BPE (Byte-Pair Encoding)}
    \begin{itemize}
        \item Fusiona iterativamente los pares de bytes (o caracteres) más frecuentes.
        \item Vocabulario de subpalabras (aprox. 50,000 tokens).
        \item Maneja palabras raras y neologismos descomponiéndolos en subpalabras conocidas.
    \end{itemize}
    % Basado en PDF y Pregunta 5
\end{frame}

\begin{frame}{Objetivo de pre-entrenamiento}
    \begin{block}{Modelado autorregresivo (causal language modeling)}
        \[
        \mathcal{L} = \sum_{t=1}^{T} \log P(x_t | x_{<t})
        \]
        Maximiza la probabilidad de cada token dados los anteriores.
    \end{block}
    \begin{itemize}
        \item No usa máscaras ni predicción de siguiente oración.
        \item El modelo aprende a predecir la siguiente palabra de forma causal.
    \end{itemize}
    % Basado en PDF y Pregunta 3
\end{frame}

\begin{frame}{Generación con GPT}
    \begin{itemize}
        \item Dado un \textbf{prompt} inicial, genera texto autoregresivamente.
        \item En cada paso, predice la siguiente palabra y la añade al contexto.
        \item Se puede usar \textbf{beam search}, \textbf{sampling} con temperatura, etc.
        \item Ideal para: generación de texto, chatbots, completar código, etc.
    \end{itemize}
    % Basado en PDF
\end{frame}

% ============================================================================
% SECCIÓN 4: COMPARACIÓN BERT VS GPT
% ============================================================================
\section{Comparación BERT vs GPT}

\begin{frame}{Tabla comparativa}
    \begin{table}
        \centering
        \begin{tabular}{l|c|c}
            \toprule
            \textbf{Característica} & \textbf{BERT} & \textbf{GPT} \\
            \midrule
            Arquitectura & Encoder-only & Decoder-only \\
            Direccionalidad & Bidireccional & Unidireccional (causal) \\
            Objetivo de entrenamiento & MLM + NSP & Modelado autorregresivo \\
            Tokenización & WordPiece & BPE \\
            Uso principal & Comprensión, clasificación, NER, QA & Generación, diálogo, completado \\
            \bottomrule
        \end{tabular}
    \end{table}
    % Basado en PDF y Pregunta 1
\end{frame}

% ============================================================================
% SECCIÓN 5: FINE-TUNING
% ============================================================================
\section{Fine-tuning}

\begin{frame}{Fine-tuning vs Feature-based}
    \begin{columns}[t]
        \column{0.5\textwidth}
        \textbf{Feature-based}
        \begin{itemize}
            \item Extraer embeddings del modelo pre-entrenado como características fijas.
            \item Entrenar un clasificador simple (SVM, MLP) sobre ellas.
            \item Más rápido, menor riesgo de sobreajuste.
        \end{itemize}
        \column{0.5\textwidth}
        \textbf{Fine-tuning}
        \begin{itemize}
            \item Añadir una cabeza específica para la tarea.
            \item Entrenar todos los parámetros (o algunos) conjuntamente.
            \item Mejor rendimiento si hay suficientes datos.
        \end{itemize}
    \end{columns}
    % Basado en PDF y Pregunta 4
\end{frame}

\begin{frame}{Tipos de fine-tuning}
    \begin{itemize}
        \item \textbf{Completo}: se actualizan todos los parámetros. Requiere más memoria y datos.
        \item \textbf{Capas congeladas}: se congelan las primeras capas (generales) y se ajustan las últimas y la cabeza. Útil con datos pequeños.
        \item \textbf{Adaptadores}: se insertan pequeñas capas entrenables entre las capas pre-entrenadas.
        \item \textbf{LoRA}: matrices de bajo rango añadidas a las matrices de pesos.
        \item \textbf{Prompt Tuning}: se aprenden ``soft prompts'' continuos; el modelo permanece congelado.
    \end{itemize}
    % Basado en PDF y Pregunta 9
\end{frame}

% ============================================================================
% SECCIÓN 6: TÉCNICAS AVANZADAS DE FINE-TUNING
% ============================================================================
\section{Técnicas Avanzadas de Fine-tuning}

\begin{frame}{LoRA (Low-Rank Adaptation)}
    \begin{itemize}
        \item Congela los pesos pre-entrenados.
        \item Añade matrices de descomposición de bajo rango a las matrices de atención y feed-forward.
        \item Para una matriz $W_0 \in \mathbb{R}^{d \times k}$, aprende $B \in \mathbb{R}^{d \times r}$, $A \in \mathbb{R}^{r \times k}$ con $r \ll \min(d,k)$.
        \item Nueva matriz: $W = W_0 + BA$.
        \item Solo se entrenan $B$ y $A$ (menos del 1\% de los parámetros).
    \end{itemize}
    % Basado en PDF y Pregunta 8
\end{frame}

\begin{frame}{Adapters y Prompt Tuning}
    \begin{columns}[t]
        \column{0.5\textwidth}
        \textbf{Adapters}
        \begin{itemize}
            \item Pequeños módulos (dos capas lineales con no linealidad) insertados entre capas del transformer.
            \item Solo se entrenan los adaptadores.
        \end{itemize}
        \column{0.5\textwidth}
        \textbf{Prompt Tuning}
        \begin{itemize}
            \item Se aprenden vectores de ``soft prompts'' que se anteponen a la entrada.
            \item El modelo base permanece congelado.
            \item Muy eficiente, funciona bien con modelos grandes.
        \end{itemize}
    \end{columns}
    % Basado en PDF y Pregunta 9
\end{frame}

% ============================================================================
% SECCIÓN 7: TOKENIZADORES DE SUBPALABRAS
% ============================================================================
\section{Tokenizadores de Subpalabras}

\begin{frame}{WordPiece, BPE y SentencePiece}
    \begin{table}
        \centering
        \begin{tabular}{l|c|c}
            \toprule
            \textbf{Algoritmo} & \textbf{Usado en} & \textbf{Métrica de fusión} \\
            \midrule
            WordPiece & BERT & $\frac{\text{freq}(ab)}{\text{freq}(a)\cdot\text{freq}(b)}$ \\
            BPE & GPT & Frecuencia bruta del par \\
            SentencePiece & T5, XLNet & BPE o Unigram (sobre caracteres Unicode, sin pre-tokenización) \\
            \bottomrule
        \end{tabular}
    \end{table}
    % Basado en PDF y Pregunta 5
\end{frame}

% ============================================================================
% SECCIÓN 8: MÉTRICAS DE EVALUACIÓN
% ============================================================================
\section{Métricas de Evaluación}

\begin{frame}{Para clasificación}
    \begin{itemize}
        \item \textbf{Accuracy}: porcentaje de predicciones correctas.
        \item \textbf{F1-score}: media armónica de precisión y recall, útil para clases desbalanceadas.
        \item \textbf{AUC-ROC}: capacidad discriminativa.
    \end{itemize}
\end{frame}

\begin{frame}{Para generación}
    \begin{itemize}
        \item \textbf{Perplejidad}: $\exp\left(-\frac{1}{N}\sum \log P(x_t|x_{<t})\right)$. Baja perplejidad = mejor modelo.
        \item \textbf{BLEU}: para traducción (precisión de n-gramas con penalización por longitud).
        \item \textbf{ROUGE}: para summarization (solapamiento de n-gramas).
        \item \textbf{BERTScore}: basada en embeddings contextuales.
    \end{itemize}
    % Basado en Pregunta 12
\end{frame}

% ============================================================================
% SECCIÓN 9: EJEMPLO PRÁCTICO - FINE-TUNING DE BERT CON HUGGING FACE
% ============================================================================
\section{Ejemplo: Fine-tuning de BERT con Hugging Face}

\begin{frame}[fragile]{Paso 1-3: Instalación, carga de datos y tokenización}
    \begin{lstlisting}
# Instalación
!pip install transformers datasets evaluate

from transformers import AutoTokenizer, AutoModelForSequenceClassification, Trainer, TrainingArguments
from datasets import load_dataset
import evaluate
import numpy as np

# Cargar datos (IMDb)
dataset = load_dataset("imdb")

# Tokenización
tokenizer = AutoTokenizer.from_pretrained("bert-base-uncased")
def tokenize_function(examples):
    return tokenizer(examples["text"], padding="max_length", truncation=True, max_length=512)
tokenized_datasets = dataset.map(tokenize_function, batched=True)
    \end{lstlisting}
    % Basado en PDF
\end{frame}

\begin{frame}[fragile]{Paso 4-6: Modelo, métrica y configuración}
    \begin{lstlisting}
# Modelo pre-entrenado
model = AutoModelForSequenceClassification.from_pretrained("bert-base-uncased", num_labels=2)

# Métrica
metric = evaluate.load("accuracy")
def compute_metrics(eval_pred):
    logits, labels = eval_pred
    predictions = np.argmax(logits, axis=-1)
    return metric.compute(predictions=predictions, references=labels)

# Configuración de entrenamiento
training_args = TrainingArguments(
    output_dir="./results",
    evaluation_strategy="epoch",
    save_strategy="epoch",
    learning_rate=2e-5,
    per_device_train_batch_size=8,
    per_device_eval_batch_size=8,
    num_train_epochs=3,
    weight_decay=0.01,
    load_best_model_at_end=True,
    metric_for_best_model="accuracy",
)
    \end{lstlisting}
\end{frame}

\begin{frame}[fragile]{Paso 7-8: Entrenamiento y evaluación}
    \begin{lstlisting}
# Trainer
trainer = Trainer(
    model=model,
    args=training_args,
    train_dataset=tokenized_datasets["train"].select(range(2000)),  # reducido para ejemplo
    eval_dataset=tokenized_datasets["test"].select(range(500)),
    compute_metrics=compute_metrics,
)

# Entrenar
trainer.train()

# Evaluar
trainer.evaluate()
    \end{lstlisting}
\end{frame}

% ============================================================================
% SECCIÓN 10: PREPARACIÓN PARA ENTREVISTAS
% ============================================================================
\section{Preparación para Entrevistas}

\begin{frame}{Preguntas típicas}
    \begin{itemize}
        \item ¿Cuál es la diferencia entre BERT y GPT?
        \item ¿Qué es el fine-tuning y cómo se diferencia de feature extraction?
        \item ¿Qué es LoRA y por qué es útil?
        \item ¿Por qué BERT necesita positional encoding?
        \item ¿Qué son MLM y NSP?
        \item ¿Cómo maneja BERT palabras desconocidas?
    \end{itemize}
    % Basado en PDF
\end{frame}

\begin{frame}{Preguntas avanzadas}
    \begin{itemize}
        \item ¿Qué es el problema de la máscara en MLM y cómo se soluciona?
        \item ¿Cómo funciona la attention mask en transformers?
        \item ¿Qué son los embeddings de segmento en BERT?
        \item ¿Cómo se evalúa un modelo de QA extractivo? (EM, F1)
        \item ¿Qué son las alucinaciones en modelos generativos y cómo mitigarlas?
    \end{itemize}
    % Basado en PDF y Preguntas 7, 11, 13, 18
\end{frame}

% ============================================================================
% SECCIÓN 11: CONEXIÓN CON EL RESTO DEL CURSO
% ============================================================================


% ============================================================================
% SECCIÓN 12: RESUMEN
% ============================================================================
\section{Resumen}

\begin{frame}{Lo que hemos aprendido}
    \begin{exampleblock}{Conceptos clave}
        \begin{itemize}
            \item \textbf{BERT}: encoder-only, bidireccional, pre-entrenado con MLM + NSP. Ideal para comprensión.
            \item \textbf{GPT}: decoder-only, unidireccional, pre-entrenado con modelado autorregresivo. Ideal para generación.
            \item \textbf{Fine-tuning}: adapta modelos pre-entrenados a tareas específicas, superando enfoques clásicos.
            \item \textbf{Técnicas eficientes}: LoRA, adapters, prompt tuning (pocos parámetros entrenables).
            \item \textbf{Tokenizadores}: WordPiece (BERT), BPE (GPT), SentencePiece.
            \item \textbf{Métricas}: accuracy, F1, perplejidad, BLEU, ROUGE.
            \item \textbf{Aplicaciones}: clasificación, NER, QA, generación de texto.
        \end{itemize}
    \end{exampleblock}
\end{frame}

% --- Diapositiva final ---
\begin{frame}
    \centering
    \Huge \textbf{¡Gracias!}

\end{frame}

\end{document}